%%
%% ICPP 2026 Paper — Multicast-Aware Attention Dataflow on Tenstorrent Many-Core
%% Based on ACM acmart sigconf template
%%

%% For submission / review:
\documentclass[sigconf, review, anonymous]{acmart}
%% For camera-ready:
% \documentclass[sigconf]{acmart}

%%
%% Useful packages
%%
\usepackage{algorithm}
\usepackage{algorithmic}
\usepackage{booktabs}
\usepackage{subcaption}
\usepackage{xspace}

%%
%% Custom commands
%%
\newcommand{\sysname}{\textsc{McastAttn}\xspace}
\newcommand{\ie}{\textit{i.e.}\xspace}
\newcommand{\eg}{\textit{e.g.}\xspace}
\newcommand{\etal}{\textit{et al.}\xspace}
\newcommand{\Bkv}{B_{\mathrm{kv}}}
\newcommand{\Ghead}{G_{\mathrm{head}}}
\newcommand{\Tgroup}{T_{\mathrm{group}}}

%%
%% Rights / conference metadata — update when known
%%
\setcopyright{acmlicensed}
\copyrightyear{2026}
\acmYear{2026}
\acmDOI{XXXXXXX.XXXXXXX}
\acmConference[ICPP '26]{ACM International Conference on Parallel Processing}{August 2026}{TBD}
\acmISBN{978-1-4503-XXXX-X/2026/08}

%%
%% Document
%%
\begin{document}

%% ============================================================
%%  TITLE
%% ============================================================
\title{Multicast-Aware Attention Dataflow and Auto-Tuning\\on Tenstorrent Many-Core for LLM Decode Inference}

%% ============================================================
%%  AUTHORS — replace with real info
%% ============================================================
\author{First Author}
\affiliation{%
  \institution{University / Lab}
  \city{City}
  \country{Country}}
\email{first@example.com}

\author{Second Author}
\affiliation{%
  \institution{University / Lab}
  \city{City}
  \country{Country}}
\email{second@example.com}

\renewcommand{\shortauthors}{Author et al.}

%% ============================================================
%%  ABSTRACT
%% ============================================================
\begin{abstract}
Large Language Model (LLM) inference is increasingly dominated by the
\emph{decode} phase, where each token generation step reads a large
KV cache from DRAM — making multi-head attention (MHA) severely
memory-bound. On many-core accelerators with tens to hundreds of
processing elements, a na\"ive mapping causes each core to redundantly
read KV data from DRAM, inflating total memory traffic linearly with
the number of consumer cores.

We present \sysname, a \emph{multicast-aware attention dataflow} and
accompanying \emph{auto-tuner} targeting the Tenstorrent Wormhole
many-core architecture. \sysname exploits Tensix's decoupled data
movers, circular buffers, and dual on-chip network (NoC) to read each
KV cache block from DRAM \emph{exactly once} and distribute it via
hardware multicast, while overlapping data movement with computation in
a pipelined schedule. An analytical cost model captures DRAM traffic,
NoC multicast latency, and compute utilization as closed-form functions
of workload parameters and a compact design space of \emph{head
grouping}, \emph{tile grouping}, \emph{block sizes}, and \emph{pipeline
depth}. The auto-tuner enumerates this space and outputs near-optimal
configurations without manual tuning.

Evaluation on Tenstorrent Wormhole N300 shows that \sysname reduces
per-step DRAM traffic by up to $X\times$ compared to the per-core-read
baseline, and improves decode latency by $Y$\% over a fixed-dataflow
configuration across representative LLaMA-style workloads ($L$ up to
4096, $d \in \{64,128\}$, $H$ up to 32). The auto-tuner finds
configurations within Z\% of exhaustive search in under one second.
\end{abstract}

%% ============================================================
%%  CCS CONCEPTS — generate at https://dl.acm.org/ccs
%% ============================================================
\begin{CCSXML}
<ccs2012>
 <concept>
  <concept_id>10010520.10010521.10010537.10010540</concept_id>
  <concept_desc>Computer systems organization~Multicore architectures</concept_desc>
  <concept_significance>500</concept_significance>
 </concept>
 <concept>
  <concept_id>10010520.10010575.10010755</concept_id>
  <concept_desc>Computer systems organization~Interconnection architectures</concept_desc>
  <concept_significance>300</concept_significance>
 </concept>
</ccs2012>
\end{CCSXML}

\ccsdesc[500]{Computer systems organization~Multicore architectures}
\ccsdesc[300]{Computer systems organization~Interconnection architectures}

\keywords{LLM inference, multi-head attention, decode, KV cache,
many-core, multicast, on-chip network, auto-tuning, Tenstorrent, Wormhole}

\maketitle

%% ============================================================
%%  1. INTRODUCTION
%% ============================================================
\section{Introduction}\label{sec:intro}

%% OUTLINE:
%% - 开场：LLM 推理需求爆发，decode 阶段逐 token 生成，注意力是核心瓶颈
%% - 问题：decode attention 极度 memory-bound，KV cache 读取占主导
%% - 众核机会：Tenstorrent Wormhole 等众核加速器有数百核与片上网络，
%%   但朴素映射导致每核独立读 DRAM → 冗余流量线性增长
%% - 缺口：FlatAttention 面向通用 tile / prefill / 仿真，
%%   未针对 Tensix 架构特性（解耦 data mover, CB, 双 NoC）和 decode 场景
%% - 本文方案一句话：多播感知数据流 + 解析代价模型 + 自动寻优
%% - 贡献列表（4 条，与 abstract 对应）

\textbf{TODO: Write full introduction.}

We make the following contributions:
\begin{itemize}
  \item A \emph{Tensix-specific attention dataflow} that reads Q/K/V
    from DRAM once and multicasts to consumer cores with decoupled
    move-compute overlap.
  \item A \emph{decode-phase KV cache multicast} design (producer–consumer
    graph, block size, single-/multi-card) that reduces DRAM traffic
    toward the theoretical minimum.
  \item An \emph{analytical cost model and auto-tuner} for head/tile
    grouping and pipeline layout that adapts to varying $(L, d, H, B)$
    without manual tuning.
  \item \emph{Evaluation on real Tenstorrent Wormhole hardware} showing
    significant DRAM and latency improvements over baselines.
\end{itemize}


%% ============================================================
%%  2. BACKGROUND AND RELATED WORK
%% ============================================================
\section{Background and Related Work}\label{sec:background}

%%% ----- 2.1 Multi-Head Attention and the Decode Phase -----

\subsection{Multi-Head Attention and the Decode Phase}\label{sec:bg-mha}

The Transformer architecture~\cite{vaswani2017attention} employs
multi-head attention (MHA) as its core sequence-modeling primitive.
For each of the $H$ attention heads, a query, key, and value triple
$(Q_h, K_h, V_h)$ is derived via learned linear projections, and
attention is computed as
\begin{equation}\label{eq:mha}
  \mathrm{Attn}_h
    = \mathrm{softmax}\!\Bigl(\frac{Q_h\, K_h^\top}{\sqrt{d}}\Bigr)\, V_h ,
\end{equation}
where $d$ denotes the per-head dimension and $L$ the sequence length,
giving $Q_h, K_h, V_h \in \mathbb{R}^{L \times d}$.
The outputs of all heads are concatenated and linearly projected to
form the layer output.

\paragraph{Prefill versus decode.}
Autoregressive LLM inference proceeds in two distinct
phases~\cite{pope2023efficiently}.
The \emph{prefill} (prompt) phase processes the entire input prompt in
a single forward pass, computing KV pairs for all $L$ input tokens
simultaneously.  Because it executes $\mathcal{O}(L^2 d H)$ FLOPs over
$\mathcal{O}(L d H)$ data elements, prefill is typically
\emph{compute-bound}.
The subsequent \emph{decode} phase generates one new token per step.
At step~$t$, only the new token's query
$q_h \in \mathbb{R}^{1 \times d}$ is produced, which attends over the
full KV cache of length~$L$:
\begin{equation}\label{eq:decode-attn}
  o_h = \mathrm{softmax}\!\Bigl(\frac{q_h\, K_h^\top}{\sqrt{d}}\Bigr)\, V_h ,
  \quad K_h, V_h \in \mathbb{R}^{L \times d}.
\end{equation}
This requires only $\mathcal{O}(L\,d)$ FLOPs per head but reads
$\mathcal{O}(L\,d)$ elements of both $K$ and $V$ from memory, yielding
an arithmetic intensity of $\mathcal{O}(1)$---firmly in the
\emph{memory-bound} regime.

\paragraph{KV cache traffic.}
During decode, each step must read the K and V caches for all $H$
heads from DRAM\@.  The total read volume per step is
\begin{equation}\label{eq:kv-traffic}
  D_{\mathrm{KV}} = 2\, L\, d\, H\, b_{\mathrm{elem}} ,
\end{equation}
where $b_{\mathrm{elem}}$ is the element size in bytes (\eg 2\,B for
\texttt{bfloat16}).
For a model with $H{=}32$ heads, $d{=}128$, and $L{=}4096$, this
amounts to ${\approx}\,64$\,MB per step---imposing a hard latency floor
set by available DRAM bandwidth.
Variants such as grouped-query attention
(GQA)~\cite{ainslie2023gqa} and multi-query attention
(MQA)~\cite{shazeer2019mqa} reduce the number of distinct KV heads,
lowering $D_{\mathrm{KV}}$ proportionally, yet the fundamental
memory-bound character of decode attention persists at practical
sequence lengths.

%%% ----- 2.2 Tenstorrent Wormhole Architecture -----

\subsection{Tenstorrent Wormhole Architecture}\label{sec:bg-wormhole}

Tenstorrent Wormhole~\cite{tenstorrent2024wormhole} is a many-core
accelerator designed for neural-network inference and training.
We summarize the architectural features most relevant to our
multicast-aware dataflow.

\paragraph{Chip and card overview.}
Each Wormhole chip integrates a grid of \emph{Tensix} processing cores
connected by a pair of on-chip networks.
The Wormhole N300 card used in our evaluation houses two chips with a
total of $P{=}128$ usable Tensix cores and 12\,GB GDDR6 DRAM\@.
Table~\ref{tab:hw-params} lists the key hardware parameters.

\begin{table}[t]
  \caption{Key hardware parameters of the Tenstorrent Wormhole N300.}
  \label{tab:hw-params}
  \centering
  \small
  \begin{tabular}{ll}
    \toprule
    \textbf{Parameter} & \textbf{Value} \\
    \midrule
    Tensix cores (N300)        & 128 (2 chips $\times$ 64 cores) \\
    L1 SRAM per core           & 1\,MB scratchpad \\
    DRAM capacity              & 12\,GB GDDR6 \\
    DRAM bandwidth (aggregate) & ${\sim}$200\,GB/s \\
    On-chip network            & Dual 2-D mesh (NoC\,0 \& NoC\,1) \\
    Multicast support          & HW single-source rectangular \\
    Compute per core           & BF16/FP16 matrix \& element-wise FMA \\
    \bottomrule
  \end{tabular}
\end{table}

\paragraph{Tensix core micro-architecture.}
Each Tensix core contains five lightweight RISC-V processors
(\emph{baby cores}) with dedicated roles:
\begin{itemize}
  \item \textbf{Two data movers} (BRISC and NCRISC) manage DMA
    transfers between DRAM, NoC, and local L1 SRAM\@.
  \item An \textbf{unpack} engine reads tiles from L1 into the compute
    pipeline's source registers.
  \item A \textbf{math/FPU} engine executes matrix-multiply and
    element-wise operations.
  \item A \textbf{pack} engine writes results from the compute pipeline
    back to L1.
\end{itemize}
Because data movers, unpack, math, and pack operate as
\emph{decoupled} pipeline stages, data transfer and computation can
overlap---a property central to the pipelined schedule described in
Section~\ref{sec:df-pipeline}.

\paragraph{Circular buffers.}
The stages within a Tensix core communicate through software-managed
\emph{circular buffers} (CBs) allocated in L1.
A producer (\eg a data mover receiving tiles from the NoC) pushes
tiles into CB slots, and a consumer (\eg the unpack engine) pops them.
By configuring the CB depth ($\geq 2$ slots), double-buffering or
deeper pipelining is achieved transparently, enabling data movement
to proceed concurrently with computation.

\paragraph{Dual NoC and hardware multicast.}
The two independent networks, NoC\,0 and NoC\,1, each span the full
2-D mesh and support both unicast and \emph{hardware multicast}.
In a multicast, a source core issues a single DMA command specifying
a rectangular region of destinations; the network replicates the
packet at branch points and delivers data to all targets in one
traversal without the source transmitting separate copies.
This multicast primitive is the \emph{key enabler} of our design:
a producer core reads a KV-cache block from DRAM once and multicasts
it to all consumer cores in its group, reducing total DRAM traffic
toward the theoretical minimum.
We assign NoC\,0 to KV distribution and NoC\,1 to partial-result
reduction, fully utilizing both networks concurrently.

%%% ----- 2.3 Related Work -----

\subsection{Related Work}\label{sec:bg-related}

\paragraph{Efficient attention on GPUs.}
FlashAttention~\cite{dao2022flashattention} introduced IO-aware tiling
to keep intermediate attention matrices in on-chip SRAM, dramatically
reducing HBM traffic on NVIDIA GPUs.
FlashAttention-2~\cite{dao2023flashattention2} improved parallelism
and work partitioning across warps and thread blocks, while
FlashAttention-3~\cite{shah2024flashattention3} further exploits
Hopper-class asynchronous execution and low-precision accumulation.
For the decode phase specifically,
Flash-Decoding~\cite{flashdecoding2023} parallelizes across the KV
sequence dimension rather than the batch or head dimension, improving
GPU utilization when the batch size is small;
FlashDecoding++~\cite{hong2024flashdecoding} extends this with a
unified max-value technique that avoids synchronization overhead in
the softmax reduction.
Ring Attention~\cite{liu2023ringattention} distributes long sequences
across multiple devices in a ring topology.
All these techniques are deeply tied to the GPU execution model
(warps, shared memory, tensor cores) and do not directly transfer to
many-core architectures with fundamentally different memory hierarchies
and communication primitives.

\paragraph{KV cache management and attention variants.}
PagedAttention~\cite{kwon2023pagedattention}, implemented in the vLLM
serving system, addresses KV cache \emph{memory fragmentation} by
borrowing virtual-memory paging concepts, enabling flexible cache
sharing across requests.
On the model-architecture side, GQA~\cite{ainslie2023gqa} and
MQA~\cite{shazeer2019mqa} reduce the number of KV heads, directly
shrinking cache footprint and per-step DRAM traffic.
StreamingLLM~\cite{xiao2024streamingllm} and
H\textsubscript{2}O~\cite{zhang2024h2o} selectively evict or compress
KV cache entries to handle long contexts.
These approaches are \emph{orthogonal} to our dataflow design: they
reduce the logical volume of KV data, whereas \sysname minimizes the
physical DRAM reads for a given cache volume via multicast.

\paragraph{LLM serving systems.}
Orca~\cite{yu2022orca} pioneered iteration-level scheduling with
continuous batching to improve GPU throughput.
DeepSpeed-Inference~\cite{aminabadi2022deepspeed} combines model
parallelism, customized kernels, and heterogeneous execution for
large-scale serving.
More recently, Splitwise~\cite{patel2024splitwise} and
DistServe~\cite{zhong2024distserve} propose \emph{disaggregating} the
prefill and decode phases onto separate hardware pools, recognizing
their distinct compute-vs-memory characteristics.
These systems orchestrate resources at the \emph{cluster or model}
level; our work is complementary, targeting kernel-level dataflow on a
single accelerator card.

\paragraph{Dataflow mapping and cost modeling.}
A rich body of work formalizes dataflow mapping for DNN accelerators.
MAESTRO~\cite{kwon2019maestro} provides an analytical framework for
evaluating data reuse, bandwidth, and energy across mappings.
Timeloop~\cite{parashar2019timeloop} and
Eyeriss~v2~\cite{chen2019eyerissv2} similarly explore the mapping
space for spatial architectures, primarily targeting convolution.
FLAT~\cite{kao2023flat} extends dataflow analysis to the attention
operator, proposing optimized mappings that mitigate the
memory-bandwidth bottleneck of softmax.
However, these tools target either generic spatial arrays or GPU-like
architectures; none models the Tensix-specific features of decoupled
data movers, per-core circular buffers, or hardware NoC multicast that
are central to our design.

\paragraph{Attention on many-core and non-GPU architectures.}
The most closely related work is
FlatAttention~\cite{flatattention2024}, which co-optimizes dataflow
and fabric collectives for multi-head attention on tile-based many-PE
accelerators.
FlatAttention demonstrates significant speedups by orchestrating data
movement across PEs via collective operations.
However, it (i)~primarily targets the \emph{prefill} phase, where the
workload is compute-bound and the optimal strategies differ from
memory-bound decode;
(ii)~employs a generic tile-PE abstraction without exploiting
architecture-specific features such as Tensix's decoupled data movers,
software-managed circular buffers, and dual-NoC hardware multicast;
and (iii)~does not provide an auto-tuner that adapts to varying decode
workload parameters.
Cerebras~\cite{cerebras2024waferscale} and RISC-V kilo-core
clusters~\cite{dynamic-alloc2024} also present many-core substrates
for transformer workloads, but focus on wafer-scale integration and
shared-memory management respectively, without addressing
decode-specific multicast dataflows.
Prior work on the Tenstorrent platform~\cite{tt-fft2025} has explored
FFT acceleration and validated the effectiveness of Wormhole's
multicast primitive for signal-processing kernels, but did not target
attention workloads.

\paragraph{Positioning of this work.}
To the best of our knowledge, no prior work jointly addresses
(i)~Tensix-specific dataflow with decoupled data movers and dual-NoC
multicast,
(ii)~the decode-phase KV-cache access pattern, and
(iii)~automated configuration tuning driven by an analytical cost
model.
\sysname fills this gap.


%% ============================================================
%%  3. PROBLEM FORMULATION
%% ============================================================
\section{Problem Formulation}\label{sec:problem}

%% OUTLINE:
%% 3.1 负载参数定义
%%   - Workload = (L, d, H, B, elem)：序列长、head 维度、head 数、batch、元素字节
%%   - 解码步计算量与访存量分析
%%
%% 3.2 优化目标
%%   - 主目标：最小化单步延迟 latency
%%   - 次目标：最小化 DRAM 流量、最大化利用率
%%   - 约束：SRAM 容量、核数、NoC 带宽
%%
%% 3.3 设计空间概览
%%   - 分组维度：G_head, T_group
%%   - 块大小：B_kv
%%   - 流水参数：depth, overlap
%%   - 归约策略：binary_tree, row_col
%%   - 空间规模 ~O(百) 种合法配置

\textbf{TODO: Formalize workload, objectives, constraints, design space.}


%% ============================================================
%%  4. MULTICAST-AWARE ATTENTION DATAFLOW
%% ============================================================
\section{Multicast-Aware Attention Dataflow}\label{sec:dataflow}

%% OUTLINE:
%% 4.1 总体架构
%%   - 核角色划分：Producer（DRAM 读 + NoC 多播）、Consumer（计算）、Reducer
%%   - 数据流示意图（Figure 1）：DRAM → Producer → multicast → Consumers → reduce → Output
%%
%% 4.2 生产者-消费者多播设计
%%   - Producer 从 DRAM 按 B_kv 块读 K/V 到本地 CB
%%   - 通过 NoC 0 单源多播到同组 T_group 个消费者核
%%   - 消费者接收到 CB 后本地计算 QK^T 分块、partial softmax、AttnV 分块
%%   - 双 NoC 利用：NoC 0 做 K/V 多播，NoC 1 做 partial result reduce
%%
%% 4.3 流水重叠调度
%%   - Block-level pipeline: DRAM read → multicast → compute
%%   - depth=1（串行）、depth=2（双缓冲）的时序图（Figure 2）
%%   - overlap=1 时多播与上一块计算并行
%%   - 稳态吞吐公式：t_steady = max(t_dram, t_mcast, t_comp)
%%
%% 4.4 片上归约
%%   - Softmax 在线归约：每块计算 local max/sum，最终跨核 reduce
%%   - binary tree reduce vs row-col reduce 策略
%%   - 输出 O 归约后由单核写回 DRAM

\subsection{Overall Architecture}\label{sec:df-arch}
\textbf{TODO: Producer/consumer/reducer role assignment, data flow
diagram (Figure~\ref{fig:dataflow}).}

\begin{figure}[t]
  \centering
  %% \includegraphics[width=\columnwidth]{figures/dataflow_overview.pdf}
  \fbox{\parbox{0.9\columnwidth}{\centering\vspace{2cm}
    \textbf{Figure Placeholder:} Dataflow overview\\
    DRAM $\to$ Producer $\to$ Multicast $\to$ Consumers $\to$ Reduce $\to$ Output
  \vspace{2cm}}}
  \caption{Overview of the \sysname multicast-aware attention dataflow
    on Tenstorrent Wormhole.}
  \label{fig:dataflow}
\end{figure}

\subsection{Producer–Consumer Multicast}\label{sec:df-mcast}
\textbf{TODO: Detail KV block multicast via NoC, CB usage, dual-NoC
assignment.}

\subsection{Pipelined Overlap Scheduling}\label{sec:df-pipeline}
\textbf{TODO: Block-level pipeline timing diagram (Figure~\ref{fig:pipeline}),
depth/overlap modes, steady-state throughput.}

\begin{figure}[t]
  \centering
  %% \includegraphics[width=\columnwidth]{figures/pipeline_schedule.pdf}
  \fbox{\parbox{0.9\columnwidth}{\centering\vspace{2cm}
    \textbf{Figure Placeholder:} Pipeline schedule\\
    Block 0: [DRAM] [mcast] [comp]\\
    Block 1: \hspace{2em} [DRAM] [mcast] [comp]\\
    Block 2: \hspace{4em} [DRAM] [mcast] [comp]
  \vspace{1cm}}}
  \caption{Block-level pipelined schedule with double buffering
    ($\text{depth}=2$, $\text{overlap}=1$).}
  \label{fig:pipeline}
\end{figure}

\subsection{On-Chip Reduction}\label{sec:df-reduce}
\textbf{TODO: Online softmax reduce, binary tree vs row-col, output write-back.}


%% ============================================================
%%  5. ANALYTICAL COST MODEL
%% ============================================================
\section{Analytical Cost Model}\label{sec:costmodel}

%% OUTLINE:
%% 5.1 DRAM 流量建模
%%   - 基线（每核读）：DRAM_bytes = n_consumers × (K_bytes + V_bytes)
%%   - 多播理想：DRAM_bytes = 1 × (K_bytes + V_bytes) + Q_bytes + O_bytes
%%   - 实际：考虑 n_groups = H / G_head 组，每组有独立 producer
%%
%% 5.2 延迟建模
%%   - 每块三阶段耗时：
%%     t_dram = B_kv × d × 2 × elem / BW_dram
%%     t_mcast = B_kv × d × 2 × elem / BW_noc (× hop factor)
%%     t_comp = B_kv × d × FLOPs_factor / FLOPS_peak
%%   - 流水稳态：t_steady 根据 (depth, overlap) 取不同 max 组合
%%   - 总延迟 = t_first_block + (n_blocks - 1) × t_steady + t_reduce + t_qo
%%
%% 5.3 利用率
%%   - utilization = t_comp_total / latency_total
%%   - 瓶颈标识：argmax(t_dram, t_mcast, t_comp)
%%
%% 5.4 SRAM 约束
%%   - 每核需缓存：2 × B_kv × d × G_head × elem × depth ≤ 45% × SRAM_per_core
%%   - 此约束剔除不可行配置

\subsection{DRAM Traffic Model}\label{sec:cm-dram}
\textbf{TODO: Closed-form DRAM bytes for baseline vs multicast.}

\subsection{Latency Model}\label{sec:cm-latency}
\textbf{TODO: Per-block stage latencies, pipeline steady-state formula,
total latency.}

\subsection{Utilization and Bottleneck Identification}\label{sec:cm-util}
\textbf{TODO: Utilization definition, bottleneck classification.}

\subsection{SRAM Feasibility Constraints}\label{sec:cm-sram}
\textbf{TODO: SRAM capacity check, infeasible config pruning.}


%% ============================================================
%%  6. AUTO-TUNER
%% ============================================================
\section{Auto-Tuner}\label{sec:tuner}

%% OUTLINE:
%% 6.1 设计空间定义
%%   - 参数及其候选值表 (Table 1)
%%     G_head ∈ {H 的因子}
%%     T_group ∈ {4, 8, 16, 32, 64, 128}（P 的因子且 ≥ 4）
%%     B_kv ∈ {32, 64, 128, 256, 512, 1024}
%%     depth ∈ {1, 2}
%%     overlap ∈ {0, 1}（需 depth ≥ 2）
%%     reduce_strategy ∈ {binary_tree, row_col}
%%   - 合法性约束过滤后空间规模约 300-500
%%
%% 6.2 搜索策略
%%   - 枚举所有合法配置 → 代价模型打分 → 按目标排序
%%   - 支持单目标（latency）和多目标（Pareto front: latency vs DRAM）
%%   - 搜索时间 < 1 秒（得益于闭式模型）
%%
%% 6.3 标定（可选）
%%   - 用少量硬件实测数据回归修正系数 β_dram, β_mcast, η_comp
%%   - 标定后重新寻优以贴合真实硬件

\subsection{Design Space}\label{sec:tuner-space}
\textbf{TODO: Parameter table (Table~\ref{tab:design-space}), constraint
filtering, space size.}

\begin{table}[t]
  \caption{Design space parameters and candidate values.}
  \label{tab:design-space}
  \begin{tabular}{lll}
    \toprule
    Parameter & Candidates & Constraint \\
    \midrule
    $\Ghead$ & Factors of $H$ & $H/\Ghead \times \Tgroup \leq P$ \\
    $\Tgroup$ & $\{4, 8, 16, \ldots, 128\}$ & Factor of $P$, $\geq 4$ \\
    $\Bkv$   & $\{32, 64, \ldots, 1024\}$ & SRAM capacity \\
    depth    & $\{1, 2\}$ & — \\
    overlap  & $\{0, 1\}$ & Requires depth $\geq 2$ \\
    reduce   & \{binary, row-col\} & — \\
    \bottomrule
  \end{tabular}
\end{table}

\subsection{Search Strategy}\label{sec:tuner-search}
\textbf{TODO: Enumerate-and-score, single-/multi-objective, runtime.}

\subsection{Calibration with Hardware Measurements}\label{sec:tuner-calib}
\textbf{TODO: Coefficient regression, calibration workflow.}


%% ============================================================
%%  7. EVALUATION
%% ============================================================
\section{Evaluation}\label{sec:eval}

%% OUTLINE:
%% 7.1 实验设置
%%   - 平台：Tenstorrent Wormhole N300 (128 Tensix cores, 12 GB DRAM)
%%   - 软件栈：TT-Metalium / TT-NN
%%   - 负载：decode attention, L ∈ {512, 1024, 2048, 4096},
%%     d ∈ {64, 128}, H ∈ {8, 16, 32}, B = 1, FP16
%%   - 基线：(1) Per-core DRAM read, (2) Fixed multicast (B_kv=64, 无重叠)
%%
%% 7.2 整体性能
%%   - 延迟对比柱状图 (Figure 3)：baseline vs fixed vs auto-tuned
%%   - DRAM 流量对比表 (Table 2)
%%   - 不同 (L, H) 组合下的改善百分比热力图 (Figure 4)
%%
%% 7.3 代价模型精度
%%   - 模型预测 vs 实测延迟散点图 (Figure 5)
%%   - MAPE (Mean Absolute Percentage Error) 统计
%%   - 标定前后精度对比
%%
%% 7.4 消融实验
%%   - B_kv 敏感性：B_kv 对延迟的影响曲线 (Figure 6)
%%   - T_group 敏感性：T_group 对延迟与归约开销的影响 (Figure 7)
%%   - 流水深度消融：depth=1 vs depth=2, overlap=0 vs 1
%%   - 各阶段耗时堆叠图 (Figure 8)
%%
%% 7.5 自动寻优器评估
%%   - 寻优器推荐 vs 穷举最优的差距
%%   - 推荐配置随 (L, d, H) 的变化趋势
%%   - 寻优器运行时间

\subsection{Experimental Setup}\label{sec:eval-setup}
\textbf{TODO: Hardware platform, software stack, workloads, baselines.}

\subsection{Overall Performance}\label{sec:eval-perf}
\textbf{TODO: Latency comparison (Figure~\ref{fig:latency-cmp}),
DRAM traffic table (Table~\ref{tab:dram}), improvement heatmap
(Figure~\ref{fig:heatmap}).}

\begin{figure}[t]
  \centering
  %% \includegraphics[width=\columnwidth]{figures/latency_comparison.pdf}
  \fbox{\parbox{0.9\columnwidth}{\centering\vspace{2.5cm}
    \textbf{Figure Placeholder:} Latency comparison bar chart\\
    Baseline / Fixed / Auto-tuned
  \vspace{1cm}}}
  \caption{Decode-step latency across workloads: per-core baseline,
    fixed multicast, and \sysname auto-tuned.}
  \label{fig:latency-cmp}
\end{figure}

\begin{table}[t]
  \caption{DRAM read volume per decode step (MB).}
  \label{tab:dram}
  \begin{tabular}{lrrr}
    \toprule
    Workload & Baseline & Fixed & \sysname \\
    \midrule
    $L\!=\!512,\; H\!=\!8$   & — & — & — \\
    $L\!=\!1024,\; H\!=\!8$  & — & — & — \\
    $L\!=\!2048,\; H\!=\!32$ & — & — & — \\
    $L\!=\!4096,\; H\!=\!32$ & — & — & — \\
    \bottomrule
  \end{tabular}
\end{table}

\begin{figure}[t]
  \centering
  %% \includegraphics[width=\columnwidth]{figures/improvement_heatmap.pdf}
  \fbox{\parbox{0.9\columnwidth}{\centering\vspace{2cm}
    \textbf{Figure Placeholder:} Improvement heatmap\\
    $L$ vs $H$, color = latency improvement \%
  \vspace{1.5cm}}}
  \caption{Latency improvement (\%) of \sysname over fixed baseline
    across $(L, H)$ combinations.}
  \label{fig:heatmap}
\end{figure}

\subsection{Cost Model Accuracy}\label{sec:eval-accuracy}
\textbf{TODO: Predicted vs measured scatter plot, MAPE, pre-/post-calibration.}

\subsection{Ablation Studies}\label{sec:eval-ablation}
\textbf{TODO: Block size sensitivity, group size sensitivity, pipeline
depth ablation, stage-breakdown stacked bar.}

\subsection{Auto-Tuner Quality}\label{sec:eval-tuner}
\textbf{TODO: Gap to exhaustive optimum, config trends, tuner runtime.}


%% ============================================================
%%  8. DISCUSSION
%% ============================================================
\section{Discussion}\label{sec:discussion}

%% OUTLINE:
%% - 局限性讨论
%%   - 当前仅针对 decode；prefill 扩展展望
%%   - 代价模型假设（带宽独占、无拥塞）的适用范围
%%   - 多卡扩展（Ethernet/Warp 互连）初步分析
%% - 与 GPU 方案的对比视角（FlashAttention on GPU vs 本方案 on Tensix）
%% - 对其他众核架构的可迁移性

\textbf{TODO: Limitations, multi-card outlook, comparison perspective
with GPU solutions, portability to other many-core architectures.}


%% ============================================================
%%  9. CONCLUSION
%% ============================================================
\section{Conclusion}\label{sec:conclusion}

%% OUTLINE:
%% - 总结三大贡献：Tensix 专用数据流 + KV 多播 + 自动寻优
%% - 关键数字回顾：DRAM 削减 X×, 延迟改善 Y%, 寻优 < 1s
%% - 未来方向：
%%   (1) 多卡 scale-out 支持
%%   (2) Prefill / 长序列 attention 扩展
%%   (3) 稀疏注意力 / MoE 场景适配
%%   (4) 与 TT-NN 编译器的集成

\textbf{TODO: Summarize contributions, key results, future work.}


%% ============================================================
%%  ACKNOWLEDGMENTS
%% ============================================================
\begin{acks}
\textbf{TODO: Acknowledgments.}
\end{acks}

%% ============================================================
%%  REFERENCES
%% ============================================================
\bibliographystyle{ACM-Reference-Format}
\bibliography{references}

\end{document}
